\documentclass{article} 
\usepackage{setting}
\begin{document}

\title{Research Statement} 
\author{Shuai Wei} 
\date{}
\maketitle

\section{Introduction}

My interests lie in commutative algebra, with emphasis on its connections to combinatorics. I mainly work with monomial ideals in the polynomial ring $R = \bbK[X_1,\dots,X_d]$ with $\bbK$ a field, specifically, \emph{edge ideals} of \emph{edge-weighted} graphs $G_\omega$ where $G = (V,E)$ is a finite simple graph with $V = \{v_1,\dots,v_d\}$ and a weight function $\omega: E \to \bbN$ with $\bbN = \{1,2,\dots\}$, as in the following example.

\begin{example} \label{uOmeMJuOmeMExample}
    Let $G_\omega$ be the following weighted graph with $V(G) = \{v_1,v_2,v_3\}$ and $E(G) = \{v_1v_2,v_2v_3,v_1v_3\}$ and $\omega(v_1v_2) = 1$, $\omega(v_2v_3) = 2$ and $\omega(v_1v_3) = 2$.
    \[
        \begin{tikzcd}
            & v_1 \ar[ld,dash,"1"'] \ar[rd,dash,"2"] \\
            v_2 \ar[rr,dash,"2"] && v_3
        \end{tikzcd}
    \]
    Then the weighted edge ideal~\cite{MR3055580} of $G_\omega$ is given by 
    \[I(G_\omega) = (X_1X_2,X_2^2X_3^2,X_1^2X_3^2)R.\]
\end{example}

    To study such ideals, in my dissertation I introduce the notion of ``generalized $\bbN$-weighted simplicial complex''. This is the topic of Section~\ref{section1} below where I construct a weighted version of Stanley-Reisner correspondence, and if it is ``bounded above'', I give an irredundant decompositon of the weighted Stanley-Reisner ideal in terms of the defined simplicial complex. At the end of this section, I define the Alexander dual in this setting and see how it is related to the dual of a monomial ideal, defined by E. Miller~\cite{miller1998alexander}.

    In Section~\ref{section2}, I discusss how I compute the Cohen-Macaulay type of weighted edge ideals combinatorially for certain edge-weighted graphs. \par

\section{\textsl{\large Generalized $\bbN$-weighted simplicial complex}} \label{section1}

    The Stanley-Reisner correspondence uses simplicial complexes to study square-free monomial ideals. In order to use similar techniques to study certain non-square-free monomial ideals, in my dissertation, I define a weighted version which I call a \emph{generalized $\bbN$-weighted simplicial complex}.

    A \emph{generalized $\bbN$-weighted simplicial complex} on $V$ is a non-empty collection $\Upsilon$ of sets of form $(V',\delta')$ with $V' \subseteq V$ and $\delta': V' \to \bbN$ that is closed under $\preccurlyeq$: For subsets $V_1',V_2' \subseteq V'$ and functions $\delta_1': V_1' \to \bbN$ and $\delta_2': V_2' \to \bbN$, we write $(V_1',\delta_1') \preccurlyeq (V_2',\delta_2')$ if $V_1' \subseteq V_2'$ and $\delta_1' \leq \delta_2'|_{V_2'}$. \par
    For a subset $V' \subseteq V$ and a function $\delta': V' \to \bbN$, we notationally set
    \[(V',\delta') = \{v_i^{\delta'(v_i)} \mid v_i \in V'\},\] 
    where $\delta'(v_i)$ is the \emph{weight} of $v_i$. \par
    For instance,
    \[\Upsilon = \{\{v_2^2,v_3\},\{v_2,v_3^2\},\{v_2,v_3\},\{v_2^2\},\{v_3^2\},\{v_1\},\{v_2\},\{v_3\},\emptyset\}\]
    is a generalized $\bbN$-weighted simplicial complex. \par
    Associated to $\Upsilon$, we define \emph{weighted Stanley-Reisner ideal} $J_{\Upsilon}$ by
\begin{align*}
    J_{\Upsilon} &= \bigl(\underline X^{(V',\delta')} \mathrel{\big |}V' \subseteq V, \ \delta': V' \to \bbN \text{ and }(V',\delta') \not\in \Upsilon\bigr)R \\
    &= \left(\prod_{v_i \in V'}X_i^{\delta'(v_i)} \mathrel{\big |}V' \subseteq V, \ \delta': V' \to \bbN \text{ and }(V',\delta') \not\in \Upsilon\right)R.
\end{align*}
If $\delta': V' \to \bbN$ is the constant function 1 for all $(V',\delta') \in \Upsilon$, then $\Upsilon$ becomes a usual simplicial complex $\Delta$, and this recovers the definition of the original notion of a Stanley-Reisner ideal. \par
    For example, given the generalized $\bbN$-weighted simplicial complex
    \[\Upsilon = \{\{v_2^2,v_3\},\{v_2,v_3^2\},\{v_2,v_3\},\{v_2^2\},\{v_3^2\},\{v_1\},\{v_2\},\{v_3\},\emptyset\},\]
    we have 
    \[ J_{\Upsilon} = (X_2^2X_3^2,X_1X_2,X_1X_3)R.\]
    \par Let $I \subseteq R$ be a monomial ideal~\cite{MR3839602}, for $j = 1,\dots,d$ assume $M_j$ is the maximal exponent of $X_j$ on the minimal generators of $I$. We define 
    \[\Upsilon(I) = \{(V',\delta') \text{ with }V' \subseteq V,\ \delta': V' \to \bbN \text{ and }\delta'(v_i) \leq M_i \fa v_i \in V' \mid \underline X^{(V',\delta')} \not\in I\},\]
    which is then a generalized $\bbN$-weighted simplicial complex. The next result from my dissertation is a weighted version of the Stanley-Reisner correspondence.

\begin{theorem}
    \begin{enumerate}
        \item 
            Let $I \subseteq R$ be a monomial ideal~\cite{MR3839602}, for $j = 1,\dots,d$ assume $M_j$ is the maximal exponent of $X_j$ on the minimal generators of $I$. Then
            \[J_{\Upsilon(I)} = I.\]
        \item
            Let $\Upsilon$ be a generalized $\bbN$-weighted simplicial complex such that for $i = 1,\dots,d$ we have $M_i^\Upsilon := \max\{\delta'(v_i) \mid v_i \in V' \text{ and } (V',\delta') \in \Upsilon\} < \infty$ . Then 
            \[\Upsilon(J_{\Upsilon}) = \Upsilon.\]
    \end{enumerate}
\end{theorem}

    For instance, with $J_\Upsilon = (X_2^2X_3^2,X_1X_2,X_1X_3)R$, we have 
    \begin{align*}
        \Upsilon(J_\Upsilon) &= \{(V',\delta') \text{ with }V' \subseteq V \text{ and } \delta': V' \to \bbN \mid \underline X^{(V',\delta')} \not\in J_{\Upsilon}\} \\
        &= \{\{v_2^2,v_3\},\{v_2,v_3^2\},\{v_2,v_3\},\{v_2^2\},\{v_3^2\},\{v_1\},\{v_2\},\{v_3\},\emptyset\} \\
        &= \Upsilon.
    \end{align*}

    In my dissertation, I show that, as in the classical setting, these yield an irrredundant irreducible decomposition of $J_{\Upsilon^\bfM}$:

\begin{theorem} \label{JUbfMDecompositionThm}
    Let $\Upsilon$ be a generalized $\bbN$-weighted simplicial complex such that for $i = 1,\dots,d$ we have $M_i^\Upsilon := \max\{\delta'(v_i) \mid v_i \in V' \text{ and } (V',\delta') \in \Upsilon\} < \infty$ . Then the ideal $J_{\Upsilon}$ has the following irreducible (primary) decomposition 
    \[J_{\Upsilon} = \bigcap_{\text{face }\rho \in \Upsilon}Q_\rho = \bigcap_{\rho \text{ facet}} Q_\rho,\]
    where $Q_\rho = (X^{\tau_\rho(v_i)} \mid v_i \in (V \smallsetminus {V'}) \sqcup W_\rho)R$, $W_\rho = \{v_i \in V' \mid \delta'(v_i) < M_i^{\Upsilon}\}$ and $\tau_\rho: (V \smallsetminus V') \sqcup W_\rho \to \bbN$ is given by
    \[\tau_\rho(v_i) = 
        \left\{\begin{array}{ll}
                1 & \text{if }v_i \in V \smallsetminus {V'}, \\
                \delta'(v_i) + 1 & \text{if }v_i \in W_\rho.
        \end{array}\right.
    \]
    The second intersection is irredundant.
\end{theorem}

For instance, with 
\[\Upsilon = \{\{v_2^2,v_3\},\{v_2,v_3^2\},\{v_2,v_3\},\{v_2^2\},\{v_3^2\},\{v_1\},\{v_2\},\{v_3\},\emptyset\},\]
and $M_1^{\Upsilon} = 1$, $M_2^{\Upsilon} = 2$ and $M_3^{\Upsilon} = 2$, we have 
\begin{align*}
    J_{\Upsilon} &= \bigcap_{\text{face }\rho \in \Upsilon}Q_\rho  \\
    &= Q_{\{v_2^2,v_3\}} \cap Q_{\{v_2,v_3^2\}} \cap Q_{\{v_2,v_3\}} \cap Q_{\{v_2^2\}} \cap Q_{\{v_3^2\}} \cap Q_{\{v_1\}} \cap Q_{\{v_2\}} \cap Q_{\{v_3\}} \cap Q_{\emptyset} \\
    &= (X_1,X_3^2)R \cap (X_1,X_2^2)R \cap (X_1,X_2^2,X_3^2)R \cap (X_1,X_3)R \cap (X_1,X_2)R \cap (X_2,X_3)R \\
    & \ \ \ \cap (X_1,X_3)R \cap (X_1,X_2)R \cap (X_1,X_2,X_3)R \\
    &= (X_1,X_3^2)R \cap (X_1,X_2^2)R \cap (X_2,X_3)R \\
    &= Q_{\{v_2^2,v_3\}} \cap Q_{\{v_2,v_3^2\}} \cap Q_{\{v_1\}} \\
    &= \bigcap_{\text{facet }\rho \in \Upsilon}Q_\rho.
\end{align*}

Much of our interest in this area focuses on the edge ideal $I(G)$ of a graph $G$~\cite{MR3822066} in some weighted settings. For the classical construction, this is the same as the Stanley-Reisner ideal $J_{\Delta_G}$ of the independence complex $\Delta_G$; this is the set of all independent subsets of $V$ where $W \subseteq V$ is \emph{independent} if no two vertices in $W$ are adjacent in $G$. In my dissertation I introduce and investigate the following weighted version of this notion. \par
    Let $G_\omega$ be an edge weighted graph on $V$. We say that an ordered pair $(V',\delta')$ with $V' \subseteq V$ and $\delta': V' \to \bbN$ satisfying $\delta'(v_i) \leq \max\{\omega(v_iv_j) \mid v_iv_j \in E\}$ for any $v_i \in V'$ is a \emph{weighted independent} set in $G_\omega$ if for any $v_i,v_j \in V'$, either $v_iv_j \not \in E$, or $\min\{\delta'(v_i),\delta'(v_j)\} < \omega(v_iv_j)$. Let $\Upsilon_\omega$ denote the set of weighted independent sets $(V',\delta')$ in $G_\omega$. This is the \emph{weighted independence complex} of $G_\omega$. \par
    The weighted independent sets in $G_\omega$ from Example~\ref{uOmeMJuOmeMExample} are by definition
    \[\Upsilon_\omega = \{\{v_1^2,v_3\},\{v_1,v_3^2\},\{v_1,v_3\},\{v_2^2,v_3\},\{v_2,v_3^2\},\{v_2,v_3\},\,\{v_1^2\},\{v_2^2\},\{v_3^2\},\{v_1\},\{v_2\},\{v_3\},\emptyset\}.\]

    The weighted edge ideal $I(G_\omega)$ can be decomposed in terms of weighted vertex covers of $G_\omega$ by~\cite{MR3055580}. In my dissertation I give the following decomposition result for these ideals in terms of weighted independent sets.
\begin{theorem} \label{wEIDecomInWIAndMIThm}
    We have $M_i^{\Upsilon_\omega} = \max\{\omega(v_iv_j) \mid v_iv_j \in E\}$ for $i = 1,\dots,d$. The ideal $J_{\Upsilon_\omega}$ has the following irreducible (primary) decompositions
    \[I(G_\omega) =  J_{\Upsilon_\omega} = \bigcap_{\rho \in \Upsilon_\omega} Q_\rho = \bigcap_{\text{maximal } \rho \in \Upsilon_\omega} Q_\rho.\]
    The second intersection is irredundant.
\end{theorem}
    For instance, in Example~\ref{uOmeMJuOmeMExample}, with $M_1^{\Upsilon_\omega} = 2$, $M_2^{\Upsilon_\omega} = 2$ and $M_3^{\Upsilon_\omega} = 2$, we have
\begin{align*}
    J_{\Upsilon_\omega} &= \bigcap_{\text{face }\rho \in \Upsilon_\omega}Q_\rho  \\
    &= Q_{\{v_1^2,v_3\}} \cap Q_{\{v_1,v_3^2\}} \cap Q_{\{v_1,v_3\}} \cap Q_{\{v_2^2,v_3\}} \cap Q_{\{v_2,v_3^2\}} \cap Q_{\{v_2,v_3\}} \\
    & \cap Q_{\{v_1^2\}} \cap Q_{\{v_2^2\}} \cap Q_{\{v_3^2\}} \cap Q_{\{v_1\}} \cap Q_{\{v_2\}} \cap Q_{\{v_3\}} \cap Q_{\emptyset} \\
    &= (X_2,X_3^2)R \cap (X_1^2,X_2)R \cap (X_1^2,X_2,X_3^2)R \cap (X_1,X_3^2)R \cap (X_1,X_2^2)R \cap (X_1,X_2^2,X_3^2)R \\
    & \cap (X_2,X_3)R \cap (X_1,X_3)R \cap (X_1,X_2)R \cap (X_2,X_3)R \cap (X_1,X_3)R \cap (X_1,X_2)R \cap (X_1,X_2,X_3)R \\
    &= (X_2,X_3^2) \cap (X_1^2,X_2) \cap (X_1,X_3^2)R \cap (X_1,X_2^2)R \\
    &= Q_{\{v_1^2,v_3\}} \cap Q_{\{v_1,v_3^2\}} \cap Q_{\{v_2^2,v_3\}} \cap Q_{\{v_2,v_3^2\}} \\
    &= \bigcap_{\text{facet }\rho \in \Upsilon_\omega}Q_\rho.
\end{align*}

    Let $\Upsilon$ be a generalized $\bbN$-weighted simplicial complex such that for $i = 1,\dots,d$ we have $M_i^\Upsilon < \infty$. Define the \emph{Alexander dual} $\Upsilon^\vee$ of $\Upsilon$ by
    \[\Upsilon^\vee = \{\rho := (V',\delta') \text{ with }V' \subseteq V, \ \delta': V' \to \bbN \mid ((V \smallsetminus {V'}) \sqcup W_\rho,\varphi_\rho) \not\in \Upsilon\},\]
    where $\varphi_\rho: (V \smallsetminus V') \sqcup W_\rho \to \bbN$ is defined by
    \[\varphi_\rho(v_i) = 
        \left\{\begin{array}{ll}
                M_i & \text{if }v_i \in V \smallsetminus {V'}, \\
                M_i-\delta'(v_i) & \text{if }v_i \in W_\rho.
        \end{array}\right.
    \]
    Then $\Upsilon^\vee$ is a generalied $\bbN$-weighted simplicial complex. \par
    For instance, the dual of the weighted independence complex $\Upsilon_\omega$ for Example~\ref{uOmeMJuOmeMExample} is
    \[\Upsilon_\omega^\vee = \{\{v_1,v_3^2\},\{v_2,v_3^2\},\{v_1,v_3\},\{v_1,v_2\},\{v_2,v_3\},\{v_1\},\{v_2\},\{v_3^2\},\{v_3\},\emptyset\}.\]
    Then
    \[J_{\Upsilon_\omega^\vee} = (X_1^2X_2,X_1X_2^2,X_1^2X_3,X_2^2X_3)R.\]
    \par In addition, I introduce and investigate a variation of E. Miller’s algebraic Alexander duality~\cite{miller1998alexander} in this context which is slightly better behaved in some respects. \par
    \par Let $I \subseteq R$ be a monomial ideal, for $j = 1,\dots,d$ assume $M_j$ and $m_j$ are the maximal and minimal exponent of $X_j$ on the minimal generators of $I$, respectively. Write $I = \bigcap_{i=1}^n P(V_i',\delta_i')$ as a irredundant irreducible decomposition, where $P(V_i',\delta_i') = (X_{j_i}^{\delta_i'(v_{j_i})} \mid v_{j_i} \in V_i')R$. Define the \emph{Alexander dual} $I^\vee$ of $I$ by
    \[I^\vee = (\underline X^{(V_i',\delta_i'')} \mid i = 1,\dots,n),\]
    where for $i = 1,\dots,n$: $\delta_i''(v_j) = M_j + m_j - \delta_i'(v_j)$ for any $v_j \in V_i'$ and $\underline X^{(V_i',\delta_i'')} = \prod_{v_i \in V_i'}X_i^{\delta''(v_i)}$. If all the $m_j$'s are 1 in the above definition, then it recovers the classical Alexander dual. In terms of this definition, I am working on proving that $I^{\vee\vee} = I$. \par
    In some situations, E. Miller's definition works better. So for the rest of this section I'll use his duality definition for monomial ideals. I justify the name of the previous construction by showing in the next two results from my dissertation that Alexander duality commutes with the weighted Stanley-Reisner correspondence.

\begin{theorem}
    Let $\Upsilon$ be a generalized $\bbN$-weighted simplicial complex such that for $i = 1,\dots,d$ we have $M_i^\Upsilon < \infty$. Then
    \[J_{\Upsilon}^\vee = J_{\Upsilon^\vee}.\]
\end{theorem}

For instance, consider the weighted independence complex $\Upsilon_\omega$ for Example~\ref{uOmeMJuOmeMExample}. Since the irredundant irreducible decomposition of $I(G_\omega)$ is 
\[I(G_\omega) = (X_1X_2,X_2^2X_3^2,X_1^2X_3^2) = (X_1,X_2^2)R \cap (X_1^2,X_2)R \cap (X_1,X_3^2)R \cap (X_2,X_3^2)R,\]
we have
\[J_{\Upsilon_\omega}^\vee = I(G_\omega)^\vee = (X_1^2X_2,X_1X_2^2,X_1^2X_3,X_2^2X_3)R = J_{\Upsilon_\omega^\vee}.\]

\begin{theorem}
    If $I \subseteq R$ is a monomial ideal, then 
    \[\Upsilon(I^\vee) = \Upsilon(I)^\vee.\]
\end{theorem}
    For instance, consider the weighted edge ideal $I(G_\omega)$ for Example~\ref{uOmeMJuOmeMExample}. We have by Macaulay2
    \[I(G_\omega)^\vee = (X_1^2X_2,X_1X_2^2,X_1^2X_3,X_2^2X_3).\]
Then 
\[\Upsilon(I(G_\omega)^\vee) = \{\{v_1,v_2\},\{v_1,v_3^2\},\{v_1,v_3\},\{v_2,v_3^2\},\{v_2,v_3\},\{v_3^2\},\{v_1\},\{v_2\},\{v_3\}, \emptyset\}.\]
We can show that $\Upsilon(I(G_\omega)) = \Upsilon_\omega$ generally. In this particular example, we have
    \begin{align*}
        \Upsilon(I(G_\omega))^\vee &= \Upsilon_\omega^\vee \\
        &= \{\{v_1,v_3^2\},\{v_2,v_3^2\},\{v_1,v_3\},\{v_1,v_2\},\{v_2,v_3\},\{v_1\},\{v_2\},\{v_3^2\},\{v_3\},\emptyset\} \\
        &= \Upsilon(I(G_\omega)^\vee).
    \end{align*}

\section{\textsl{\large Cohen-Macaulay type of weighted edge ideals}} \label{section2}
The \emph{type}~\cite{MR1251956} of a graded Cohen-Macaulay quotient ring $R/I$ 
    \[\operatorname{r}_R(R/I) = \dim_\bbK(\operatorname{Ext}_R^n(\bbK,R/I))~\cite{sather2009homological}\]
    quantifies important aspects of the minimal injective resolution of $R/I$ and is accordingly a measure of how close the ring is to being Gorenstein~\cite{MR0242802}. However, this invariant is generally quite difficult to compute. In my dissertation, I explicitly calculate this number combinatorially when $I$ is a monomial ideal constructed from an edge-weighted graph. \par
    For an edge-weighted graph $G_\omega$, Paulsen and Sather-Wagstaff introduce a notion of weighted vertex cover and minimal weighted vertex cover~\cite{MR3055580} that is similar to the unweighted notions, but that also tracks weighted information. For instance, the minimal weighted vertex covers for the weighted 2-path 
    \[(P_2)_\epsilon = (\begin{tikzcd}v_1 \ar[r,dash,"2"] & v_2 \ar[r,dash,"3"] & v_3)\end{tikzcd}\]
    are indicated graphically by circling the vertices in the cover and giving them weights that ``cover'' all the weighted edges of the graph:
    \[
        \begin{tikzcd}[every matrix/.append style={name=m},
          execute at end picture={
              \draw [<-] (m-1-1) ellipse (0.33cm and 0.26cm);
              \draw [<-] (m-1-3) ellipse (0.33cm and 0.26cm);
          }]
          v_1^2 \ar[r,dash,"2"] & v_2 \ar[r,dash,"3"] & v_3^3
        \end{tikzcd} \ \ \ \ \ \ 
        \begin{tikzcd}[every matrix/.append style={name=m},
          execute at end picture={
              \draw [<-] (m-1-2) ellipse (0.33cm and 0.26cm);
          }]
          v_1 \ar[r,dash,"2"] & v_2^2 \ar[r,dash,"3"] & v_3.
        \end{tikzcd}
    \]
    Paulsen and Sather-Wagstaff~\cite{MR3055580} also characterize the weighted trees that are Cohen-Macaulay. Their characterization is similar to Villarreal's description of the unweighted Cohen-Macaulay trees~\cite{MR1031197}. In particular, they are all weighted ``suspensions,'' which we write as $G_\omega=(\Sigma H)_\omega$ where $H_\epsilon$ is an induced weighted subgraph. In the following theorem in my dissertation, I show how to compute the type of $R'/I((\Sigma H)_\omega)$ combinatorially in terms of the information of $H_\epsilon$.

\begin{theorem} \label{weightedCMTypeEqualToMinWVCTheorem}
    Let $(\Sigma H)_\omega$ be a Cohen-Macaulay weighted suspension~\cite{MR3055580} of $H_\epsilon$ such that $\omega(v_iv_j) \leq \omega(v_iw_i)$ and $\omega(v_iv_j) \leq \omega(w_jv_j)$ for each $v_iv_j \in E$. Then
    \[\operatorname{r}_{R'}\biggl(\frac{R'}{I((\Sigma H)_\omega)}\biggr) = \sharp\, \{\text{minimal weighted vertex covers of }H_\epsilon\}.\]
\end{theorem}

    For example, an example of a Cohen-Macaulay weighted suspension $(\Sigma P_2)_\omega$ of 
    \[(P_2)_\epsilon = (\begin{tikzcd}v_1 \ar[r,dash,"2"] & v_2 \ar[r,dash,"3"] & v_3\end{tikzcd})\]
    is 
    \[
        \begin{tikzcd}
            w_1 \ar[d,dash,"5"] & w_2 \ar[d,dash,"3"] & w_3 \ar[d,dash,"4"] \\
            v_1 \ar[r,dash,"2"] & v_2 \ar[r,dash,"3"] & v_3.
        \end{tikzcd}
    \]
    Then since we calculated the minimal weighted vertex covers of $(P_2)_\epsilon$ above, we have
    \[\operatorname{r}_{R'}\bigl(R'/I((\Sigma P_2)_\omega)\bigr) = \sharp\, \{\text{minimal weighted vertex covers of }(P_2)_\epsilon\} = 2.\]

\section{\textsl{\large Future Work}}

I will continue to study the Cohen-Macaulayness of monomial ideals $I$ by studying $\Upsilon(I)$ combinatorially and I will also investigate when $I$ is sequentially Cohen-Macaulay and describe the minimal free resolution of $I$ using these ideas. 

Furthermore, I am interested in classifying certain other monomial ideal constructions (e.g., path ideals~\cite{MR3175038,MR3335988}) for edge-weighted graphs which are Cohen-Macaulay, and then computing the type combinatorially. \par

\bibliographystyle{plain} 
\bibliography{bibliography}

\end{document}


%%
%% End of file `elsarticle-template-num.tex'.
